The results for $C=1$ are summarized in \autoref{tab:threshold_interpretable_results_C_1}.
We report accuracy (Acc.), precision (Prec.) and recall (Rec.).
A threshold is shown as both an absolute value (i.e. $-0.20$) and a quantile level
with respect to the observed activity scores of this TF in the training dataset (i.e. $Q(0.50)$).
With merely one rule, a remarkable accuracy can be achieved.
Increasing $D$ consistently yields the same rule, plus a new clause.
Furthermore, the rulesets found for $Q_y=0.5$ and $Q_y=0.6$ are distinct.
There is a performance jump going from $D=1$ to $D=2$, less so going from $D=2$ to $D=3$.
This suggests considering a compound rule, rather than simply a clause, has some effect, but
adding more clauses in the compound rule does not appear to have much effect.

The results for $C=2$ are summarized in \autoref{tab:threshold_interpretable_results_C_2}.
By allowing for two rules, accuracy only increases marginally.
Increasing $D$ from $2$ to $3$ again yields practically no performance gain.
Notably, increasing $D$ now yields qualitatively different rulesets.
Another observation is that for $Q_y=0.5$, $C=2$, $D=2$, both rules involve the same TF\@.

\begin{table}[h]
    \centering
    \bwtable{gen/PDL1ex2/PDL1ex2_tab_interpretable_mw_0_C_1}
    \caption{Results for the full Broad dataset for $C=1$ and small rules ($D=1,2,3$).}
    \label{tab:threshold_interpretable_results_C_1}
\end{table}

\begin{table}[h]
    \bwtable{gen/PDL1ex2/PDL1ex2_tab_interpretable_mw_0_C_2}
    \caption{Results for the full Broad dataset for $C=2$ and small rules ($D=1,2,3$).}
    \label{tab:threshold_interpretable_results_C_2}
\end{table}

The results for the second experiment are shown in \autoref{fig:threshold_accuracy_mw_0.svg}.
First observe that setting $C=5, \; D=3$ in the threshold model
is comparable to setting $C=20, \; D=3$ in the boolean model.
Therefore, compared to the experiment described in \autoref{sec:boolean-methodology},
this experiment considers lower ruleset complexity.
As seen in previous tables, low values of $C$ yield reasonable performance.
While accuracy increases with $C$, the increase is marginal considering the increase in perceived ruleset complexity.
To clarify, a table like \autoref{tab:threshold_interpretable_results_C_1} for $C=5, D=3$ would span several pages.
Increasing $D$ has little effect on performance and even $D=1$ shows reasonable performance for low $Q_y$.
However, as seen in the previous experiment, there is some performance jump going from $D=1$ to $D=2$.

The cross-validation results for $Q_y=0.5, 0.6$ and $D=2$ are shown in
\autoref{fig:threshold_train_test_accuracy_mw_0.svg}.
If $C=1$, then train and test accuracy are close.
Increasing $C$ marginally improves training accuracy, but does not significantly improve test accuracy.
Hence, it appears that increasing $C$ causes overfitting rather than learning.

At this point, it appears that the most effective rulesets are the absolute simplest ones:
those with one rule and only one or two TF\@.
It may be that more complicated rulesets simply do not perform better on this dataset.
It may also be that the model does not behave as intended.
Hence, we study an alternative approach with an independent Python implementation in \autoref{sec:threshold_greedy}.

Finally, the ruleset frequency counts are given in
\autoref{tab:threshold_cross_interpretable_mw_0_C_1}.
This reveals that rulesets are sometimes (nearly) consistent, but not always.
A possible explanation for the inconsistency is that some TF are known to bind, for instance RELA and NFKB1.
It is therefore possible that the complex claimed to be promoting PD-L1 expression is a complex of
RELA, NFKB1 and STAT2.
However, NFKB1 only appears once in the entire table.
More generally, the high cross-correlation between TF activity scores as observed in \autoref{fig:corr_activity}
suggests that TF could be effectively interchangeable in the dataset.
Whether this is a real-life phenomenon or an artifact of the TF estimation technique is unclear.

\begin{figure}[h]
    \centering
    \includesvg[width=0.7\textwidth]{gen/PDL1ex2/PDL1ex2_fig_accuracy_mw_0.svg}
    \caption{Accuracy of results for the full Broad dataset for $C=1,2,3,4,5$, $D=1,2,3,4$ and $Q_y=0.5,0.6,0.7,0.8,0.9.$}
    \label{fig:threshold_accuracy_mw_0.svg}
\end{figure}

\begin{figure}[h]
    \centering
    \includesvg[width=0.8\textwidth]{gen/PDL1ex2/fig_train_test_accuracy_mw_0.svg}
    \caption{Confidence intervals ($95\%)$ for train and test accuracy of results under 10-fold stratified cross-validation on the Broad dataset, $C=1,2,3.$}
    \label{fig:threshold_train_test_accuracy_mw_0.svg}
\end{figure}

\begin{table}[h]
    \centering
    \begin{tabular}{llllr}
\toprule
 &  &  &  & \textbf{Frequency} \\
$\mathbf{Q_y}$ & $\mathbf{C}$ & $\mathbf{D}$ & \textbf{TF combination} &  \\
\midrule
\multirow[t]{6}{*}{0.5} & \multirow[t]{6}{*}{1} & \multirow[t]{3}{*}{2} & {RELA $\wedge$ STAT2} & 6 \\
 &  &  & {FOS $\wedge$ STAT2} & 3 \\
 &  &  & {NFKB1 $\wedge$ STAT2} & 1 \\
\cline{3-5}
 &  & \multirow[t]{3}{*}{3} & {FOS $\wedge$ IRF1 $\wedge$ STAT2} & 5 \\
 &  &  & {FOS $\wedge$ RELA $\wedge$ STAT2} & 3 \\
 &  &  & {IRF1 $\wedge$ RELA $\wedge$ STAT2} & 2 \\
\cline{1-5} \cline{2-5} \cline{3-5}
\multirow[t]{3}{*}{0.6} & \multirow[t]{3}{*}{1} & 2 & {FOSL1 $\wedge$ STAT2} & 10 \\
\cline{3-5}
 &  & \multirow[t]{2}{*}{3} & {FOS $\wedge$ FOSL1 $\wedge$ STAT2} & 9 \\
 &  &  & {FOSL1 $\wedge$ STAT2 $\wedge$ STAT6} & 1 \\
\cline{1-5} \cline{2-5} \cline{3-5}
\bottomrule
\end{tabular}

    \caption{Ruleset frequency count for the results of 10-fold stratified cross validation on the Broad dataset.
    In the frequency count, the threshold values and directions have been disregarded.
    Therefore, this table shows how often distinct TF combinations appear when only one rule is to be deduced from a cross-validation split of the Broad dataset}
    \label{tab:threshold_cross_interpretable_mw_0_C_1}
\end{table}
%
%\begin{table}[h]
%    \centering
%    \begin{tabular}{llllr}
\toprule
 &  &  &  & \textbf{Frequency} \\
$\mathbf{Q_y}$ & $\mathbf{C}$ & $\mathbf{D}$ & \textbf{TF combination} &  \\
\midrule
\multirow[t]{4}{*}{0.5} & \multirow[t]{4}{*}{1} & \multirow[t]{2}{*}{2} & {RELA $\wedge$ STAT2} & 7 \\
 &  &  & {FOS $\wedge$ STAT2} & 3 \\
\cline{3-5}
 &  & \multirow[t]{2}{*}{3} & {FOS $\wedge$ RELA $\wedge$ STAT2} & 9 \\
 &  &  & {FOS $\wedge$ IRF1 $\wedge$ STAT2} & 1 \\
\cline{1-5} \cline{2-5} \cline{3-5}
\multirow[t]{5}{*}{0.6} & \multirow[t]{5}{*}{1} & \multirow[t]{2}{*}{2} & {FOSL1 $\wedge$ STAT2} & 9 \\
 &  &  & {RELA $\wedge$ STAT2} & 1 \\
\cline{3-5}
 &  & \multirow[t]{3}{*}{3} & {FOS $\wedge$ FOSL1 $\wedge$ STAT2} & 8 \\
 &  &  & {FOSL1 $\wedge$ STAT2 $\wedge$ STAT6} & 1 \\
 &  &  & {FOSL1 $\wedge$ RELA $\wedge$ STAT2} & 1 \\
\cline{1-5} \cline{2-5} \cline{3-5}
\bottomrule
\end{tabular}

%    \label{tab:threshold_cross_interpretable_mw_0.2_C_1}
%\end{table}
%
%\begin{table}[h]
%    \centering
%    \begin{tabular}{llllr}
\toprule
 &  &  &  & \textbf{Frequency} \\
$\mathbf{Q_y}$ & $\mathbf{C}$ & $\mathbf{D}$ & \textbf{TF combination} &  \\
\midrule
\multirow[t]{3}{*}{0.5} & \multirow[t]{3}{*}{1} & 2 & {RELA $\wedge$ STAT2} & 10 \\
\cline{3-5}
 &  & \multirow[t]{2}{*}{3} & {FOS $\wedge$ RELA $\wedge$ STAT2} & 9 \\
 &  &  & {RELA $\wedge$ STAT2 $\wedge$ STAT6} & 1 \\
\cline{1-5} \cline{2-5} \cline{3-5}
\multirow[t]{7}{*}{0.6} & \multirow[t]{7}{*}{1} & \multirow[t]{3}{*}{2} & {FOSL1 $\wedge$ STAT2} & 4 \\
 &  &  & {FOSL1 $\wedge$ RELA} & 3 \\
 &  &  & {RELA $\wedge$ STAT2} & 3 \\
\cline{3-5}
 &  & \multirow[t]{4}{*}{3} & {FOSL1 $\wedge$ RELA $\wedge$ STAT2} & 7 \\
 &  &  & {FOSL1 $\wedge$ IRF1 $\wedge$ RELA} & 1 \\
 &  &  & {FOS $\wedge$ FOSL1 $\wedge$ STAT2} & 1 \\
 &  &  & {FOS $\wedge$ RELA $\wedge$ STAT2} & 1 \\
\cline{1-5} \cline{2-5} \cline{3-5}
\bottomrule
\end{tabular}

%    \label{tab:threshold_cross_interpretable_mw_0.4_C_1}
%\end{table}
%
%\begin{table}[h]
%    \centering
%    \begin{tabular}{llllr}
\toprule
 &  &  &  & \textbf{Frequency} \\
$\mathbf{Q_y}$ & $\mathbf{C}$ & $\mathbf{D}$ & \textbf{TF combination} &  \\
\midrule
\multirow[t]{2}{*}{0.5} & \multirow[t]{2}{*}{1} & 2 & {RELA $\wedge$ STAT2} & 10 \\
\cline{3-5}
 &  & 3 & {FOS $\wedge$ RELA $\wedge$ STAT2} & 10 \\
\cline{1-5} \cline{2-5} \cline{3-5}
\multirow[t]{6}{*}{0.6} & \multirow[t]{6}{*}{1} & \multirow[t]{2}{*}{2} & {RELA $\wedge$ STAT2} & 5 \\
 &  &  & {FOSL1 $\wedge$ RELA} & 5 \\
\cline{3-5}
 &  & \multirow[t]{4}{*}{3} & {FOS $\wedge$ RELA $\wedge$ STAT2} & 4 \\
 &  &  & {FOSL1 $\wedge$ RELA $\wedge$ STAT2} & 3 \\
 &  &  & {FOSL1 $\wedge$ IRF1 $\wedge$ RELA} & 2 \\
 &  &  & {FOS $\wedge$ FOSL1 $\wedge$ RELA} & 1 \\
\cline{1-5} \cline{2-5} \cline{3-5}
\bottomrule
\end{tabular}

%    \label{tab:threshold_cross_interpretable_mw_0.6_C_1}
%\end{table}
%
%\begin{table}[h]
%    \centering
%    \begin{tabular}{llllr}
\toprule
 &  &  &  & \textbf{Frequency} \\
$\mathbf{Q_y}$ & $\mathbf{C}$ & $\mathbf{D}$ & \textbf{TF combination} &  \\
\midrule
\multirow[t]{2}{*}{0.5} & \multirow[t]{2}{*}{1} & 2 & {RELA $\wedge$ STAT2} & 10 \\
\cline{3-5}
 &  & 3 & {FOS $\wedge$ RELA $\wedge$ STAT2} & 10 \\
\cline{1-5} \cline{2-5} \cline{3-5}
\multirow[t]{6}{*}{0.6} & \multirow[t]{6}{*}{1} & \multirow[t]{2}{*}{2} & {RELA $\wedge$ STAT2} & 6 \\
 &  &  & {FOSL1 $\wedge$ RELA} & 4 \\
\cline{3-5}
 &  & \multirow[t]{4}{*}{3} & {FOS $\wedge$ RELA $\wedge$ STAT2} & 5 \\
 &  &  & {FOSL1 $\wedge$ IRF1 $\wedge$ RELA} & 2 \\
 &  &  & {FOS $\wedge$ FOSL1 $\wedge$ RELA} & 2 \\
 &  &  & {FOSL1 $\wedge$ RELA $\wedge$ STAT2} & 1 \\
\cline{1-5} \cline{2-5} \cline{3-5}
\bottomrule
\end{tabular}

%    \label{tab:threshold_cross_interpretable_mw_0.8_C_1}
%\end{table}
%
%\begin{table}[h]
%    \centering
%    \begin{tabular}{llllr}
\toprule
 &  &  &  & \textbf{Frequency} \\
$\mathbf{Q_y}$ & $\mathbf{C}$ & $\mathbf{D}$ & \textbf{TF combination} &  \\
\midrule
\multirow[t]{2}{*}{0.5} & \multirow[t]{2}{*}{1} & 2 & {RELA $\wedge$ STAT2} & 10 \\
\cline{3-5}
 &  & 3 & {FOS $\wedge$ RELA $\wedge$ STAT2} & 10 \\
\cline{1-5} \cline{2-5} \cline{3-5}
\multirow[t]{6}{*}{0.6} & \multirow[t]{6}{*}{1} & \multirow[t]{2}{*}{2} & {RELA $\wedge$ STAT2} & 6 \\
 &  &  & {FOSL1 $\wedge$ RELA} & 4 \\
\cline{3-5}
 &  & \multirow[t]{4}{*}{3} & {FOS $\wedge$ RELA $\wedge$ STAT2} & 5 \\
 &  &  & {FOSL1 $\wedge$ IRF1 $\wedge$ RELA} & 2 \\
 &  &  & {FOS $\wedge$ FOSL1 $\wedge$ RELA} & 2 \\
 &  &  & {FOSL1 $\wedge$ RELA $\wedge$ STAT2} & 1 \\
\cline{1-5} \cline{2-5} \cline{3-5}
\bottomrule
\end{tabular}

%    \label{tab:threshold_cross_interpretable_mw_1.0_C_1}
%\end{table}
%
